% ============================================================
% Chapter 5: Development
%   Section: News Acquisition: Discovery and Extraction
%     Subsection: Discovery by Topic Group
% Included from chapters/05-development/02-news-acquisition/news-acquisition.tex
% ============================================================

\subsection{Discovery by Topic Group}
\label{subsec:scraping-groups}

Until 5~October~2026, ingestion read every enabled source and sent
every new article it found to a full analysis: a scrape, an enrichment
and one language-model call per claim, for each. Discovery is now
scoped to what the editor asks for.

\paragraph{Five groups.} The 23 topics of the topic classifier are
organised into five groups: society, science, environment, culture and
health, the same groups over which the custom validation set is
balanced (Section~\ref{sec:custom-dataset}). They are defined once, in
\path{src/config/topics.py}. A round asks for groups rather than topics
because several topics are covered by only one or two sources, so a
topic picker would often return nothing; and it asks for at most three,
because four of five is close to everything, which is what scoping
exists to avoid.

\paragraph{Which sources a group reads.} Each source file lists its
\texttt{groups}: a decision taken by hand from what the source actually
publishes, not derived from its free-form tags, so a generalist outlet
names only the groups it has real sections for. Tests hold every
enabled source to at least one group, and every group to at least one
English and one Spanish source, so no choice of groups returns a list
in one language only; an unknown group name is refused when the file
is read. Table~\ref{tab:group-coverage} shows the coverage of the
enabled sources.

\begin{table}[ht]
\centering
\small
\caption{Enabled sources per topic group and language. A source may
cover several groups.}
\label{tab:group-coverage}
\begin{tabular}{@{}lrr@{}}
\toprule
Group & English & Spanish \\
\midrule
Society     &  8 & 9 \\
Science     & 11 & 11 \\
Environment & 11 & 3 \\
Culture     &  4 & 5 \\
Health      &  8 & 7 \\
\bottomrule
\end{tabular}
\end{table}

Only the sources that name a requested group are read, and only for
the topics of the requested groups: the section names searched, the
keywords matched and the links accepted all follow from those topics.

\paragraph{What it saves.} The discovery experiment
(\texttt{python -m src.evaluation.experiments discovery}) measures each
setting against a baseline that reads every source for every topic. On
5~October~2026, with two candidates per source, Environment alone read
14 of the 34 sources then enabled ($0.41\times$) and 153 of 1,536 links
($0.10\times$), in 10~s against 15~s; Culture and Health together read
23 sources ($0.68\times$) and 596 links ($0.39\times$). Feeds change
during the day, so settings are compared within one run, and repeated
on several days for an interval (Section~\ref{sec:experiments-pre-verdict}).

The table also shows the catalogue's weak spot: Environment has three
Spanish sources against eleven English ones, so an Environment round
is mostly in English.
